\chapter{Unified Steppable Solver — Phase C}
\label{ch:steppable_solver_concept}

% =============================================================================
%  Chapter 61: SteppableSolver concept, NLAnalysis single-step API,
%              and AnalysisState cross-engine transfer
% =============================================================================

\section{Motivation}

Chapter~59 introduced \texttt{IncrementalControlPolicy} for injectable
load-stepping strategies in \texttt{NonlinearAnalysis}, while Chapter~60
added single-step control (\texttt{step()}, \texttt{step\_to()},
\texttt{step\_n()}) and condition-based breakpoints (\texttt{StepDirector})
to \texttt{DynamicAnalysis}.

Phase~C closes the gap by:
\begin{enumerate}
  \item Defining a C++23 \textbf{concept} \texttt{SteppableSolver} that
        both engines satisfy.
  \item Extending \texttt{NonlinearAnalysis} with the same single-step
        API already available in \texttt{DynamicAnalysis}.
  \item Providing \texttt{AnalysisState}, a lightweight snapshot struct
        for cross-engine state transfer.
\end{enumerate}

The result is a unified stepping protocol: any algorithm templated on
\texttt{SteppableSolver} works identically with static or dynamic solvers.


\section{The \texttt{SteppableSolver} Concept}

\begin{verbatim}
namespace fall_n {

template <typename S>
concept SteppableSolver = requires(S& solver, const S& csolver,
                                   double target, int n)
{
    { solver.step()          } -> std::convertible_to<bool>;
    { solver.step_to(target) } -> std::same_as<StepVerdict>;
    { solver.step_n(n)       } -> std::same_as<StepVerdict>;
    { csolver.current_time() } -> std::convertible_to<double>;
    { csolver.current_step() } -> std::convertible_to<PetscInt>;
};

} // namespace fall_n
\end{verbatim}

The ``time axis'' has engine-specific semantics:
\begin{itemize}
  \item \textbf{DynamicAnalysis}: physical time $t$ (seconds).
  \item \textbf{NonlinearAnalysis}: control parameter $p \in [0, 1]$
        (0 = unloaded, 1 = full load/displacement).
\end{itemize}

Both engines advance monotonically along their axis and share the same
\texttt{StepVerdict} / \texttt{StepDirector} infrastructure from
Chapter~60.


\section{NonlinearAnalysis Single-Step API}

\subsection{Initialization}

Before calling \texttt{step()}, the user must initialize the stepping
state with \texttt{begin\_incremental()}:

\begin{verbatim}
NonlinearAnalysis<Policy> nl{&model};

// Default load control, 10 steps, max 4 bisection levels
nl.begin_incremental(10);

// Or with an explicit control scheme
nl.begin_incremental(10, 4, DisplacementControl{1, 0, 0.01});
\end{verbatim}

Internally, \texttt{begin\_incremental()} calls \texttt{setup()}, zeroes
the solution vector, captures the reference force vector~$\mathbf{f}_\text{full}$,
computes the increment size $\Delta p = 1/N$, and type-erases the control
scheme into a \texttt{std::function} for subsequent calls.

\subsection{Stepping}

\begin{verbatim}
bool ok = nl.step();            // advance one increment (dp)
auto v  = nl.step_to(0.5);     // advance to p = 0.5
auto v  = nl.step_n(3);        // advance 3 increments
\end{verbatim}

Each \texttt{step()} call:
\begin{enumerate}
  \item Computes $p_\text{target} = p_\text{done} + \Delta p$.
  \item Calls the bisection engine
        \texttt{advance\_step\_impl\_(p\_done, p\_target, max\_bisections)}.
  \item On convergence: updates $p_\text{done}$, increments the step counter,
        fires the step callback, and calls
        \texttt{model\_.update\_elements\_state()}.
  \item Returns \texttt{true} on convergence, \texttt{false} on divergence.
\end{enumerate}

\subsection{Runtime Reconfiguration}

\begin{verbatim}
nl.set_increment_size(0.05);    // change dp at runtime
nl.set_director(my_director);   // install a StepDirector
\end{verbatim}

The \texttt{StepDirector} factories from Chapter~60 work identically:
\texttt{pause\_at\_times}, \texttt{pause\_every\_n}, \texttt{pause\_on},
\texttt{stop\_on}, and \texttt{compose}.

In the \texttt{StepEvent} passed to the director:
\begin{itemize}
  \item \texttt{step} = load step count.
  \item \texttt{time} = control parameter $p$.
  \item \texttt{displacement} = global displacement vector $\mathbf{U}$.
  \item \texttt{velocity} = \texttt{nullptr} (no velocity in static analysis).
\end{itemize}


\section{Type-Erased Bisection Engine}

The original \texttt{advance\_to\_p\_<CS>()} is templated on the control
scheme type.  For the single-step API, where the scheme type is erased at
\texttt{begin\_incremental()} time, we add a non-template companion
\texttt{advance\_step\_impl\_()} that uses the stored
\texttt{std::function<void(double, Vec, Vec, ModelT*)>}.

The algorithm is identical (snapshot, apply, solve, commit-or-bisect) but
operates on the member \texttt{apply\_fn\_} and \texttt{f\_full\_} instead
of template parameters.  The original templated version is preserved for
backward compatibility of \texttt{solve\_incremental()}.


\section{AnalysisState — Cross-Engine Transfer}

\begin{verbatim}
namespace fall_n {

struct AnalysisState {
    Vec      displacement{nullptr};
    Vec      velocity{nullptr};       // nullptr for static
    double   time{0.0};
    PetscInt step{0};
};

} // namespace fall_n
\end{verbatim}

Both \texttt{NonlinearAnalysis} and \texttt{DynamicAnalysis} provide
\texttt{get\_analysis\_state()} returning a snapshot with borrowed PETSc
vectors.  Typical usage for a preload $\to$ dynamic handoff:

\begin{verbatim}
// Phase 1: static preload
NonlinearAnalysis<Policy> nl{&model};
nl.begin_incremental(10);
nl.step_to(1.0);

auto state = nl.get_analysis_state();

// Phase 2: dynamic excitation (transfer displacement)
DynamicAnalysis<Policy> dyn{&model};
dyn.set_initial_displacement(state.displacement);
dyn.step_to(10.0);
\end{verbatim}


\section{Verification}

\begin{center}
\begin{tabular}{clr}
\hline
\textbf{\#} & \textbf{Test} & \textbf{Assertions} \\
\hline
1  & \texttt{SteppableSolver} concept satisfaction (NL + Dynamic)  & 2  \\
2  & NL \texttt{begin\_incremental + step()}                       & 6  \\
3  & NL \texttt{step\_n(n)}                                        & 6  \\
4  & NL \texttt{step\_to(p)}                                       & 6  \\
5  & NL \texttt{step\_to + pause\_at\_times}                       & 6  \\
6  & NL \texttt{step\_to + pause\_on} (custom condition)           & 2  \\
7  & NL \texttt{set\_increment\_size} (runtime dp change)          & 4  \\
8  & NL equivalence (step-by-step vs \texttt{solve\_incremental})  & 7  \\
9  & NL \texttt{get\_analysis\_state}                              & 5  \\
10 & Dynamic \texttt{get\_analysis\_state}                         & 4  \\
\hline
   & \textbf{Total}                                                & \textbf{48} \\
\hline
\end{tabular}
\end{center}

Test~8 confirms exact equivalence: the step-by-step path produces
$\|\Delta\mathbf{u}\| = 0$ compared to the batch
\texttt{solve\_incremental()} path.

All 47/51 suite tests pass (4 pre-existing data-file failures unchanged).


\section{Files}

\begin{description}
  \item[\texttt{src/analysis/SteppableSolver.hh}] \hfill \\
    \texttt{SteppableSolver} concept + \texttt{AnalysisState} struct.
  \item[\texttt{src/analysis/NLAnalysis.hh}] \hfill \\
    Added: stepping state members, \texttt{begin\_incremental()},
    \texttt{step()}, \texttt{step\_to()}, \texttt{step\_n()},
    \texttt{current\_time()}, \texttt{current\_step()},
    \texttt{set\_director()}, \texttt{set\_increment\_size()},
    \texttt{get\_analysis\_state()}, \texttt{advance\_step\_impl\_()}.
  \item[\texttt{src/analysis/DynamicAnalysis.hh}] \hfill \\
    Added: \texttt{get\_analysis\_state()}, \texttt{SteppableSolver.hh} include.
  \item[\texttt{src/analysis/Analysis.hh}] \hfill \\
    Added: \texttt{SteppableSolver.hh} include.
  \item[\texttt{tests/test\_steppable\_solver.cpp}] \hfill \\
    10 tests, 48 assertions.
  \item[\texttt{CMakeLists.txt}] \hfill \\
    Registered \texttt{fall\_n\_steppable\_solver\_test}.
\end{description}
